\section{Training Interventions and Reproduction}
\label{app:training}

This appendix specifies the optimized objectives, the parameters reached by each gradient, and the student state available after training.
Appendix~\ref{app:data-targets} defines the data, targets, and controls; Appendix~\ref{app:inference} defines the fresh readouts and statistical inference; and Section~\ref{sec:results} with Appendices~\ref{app:secondary}--\ref{app:sensitivities} report the resulting observations and verdicts.
Temporary target heads never enter a later evaluation.

\subsection{Ordinary \ac{kd} and language-quality rules}

Let \(z_T(x)_t\) and \(z_S(x)_t\) be teacher and student logits at a non-padding causal-language-model position \(t\).
At temperature \(\tau\), their softened token distributions are
\[
  p_T^\tau(\cdot\mid x,t)
  =\operatorname{softmax}\!\left(z_T(x)_t/\tau\right),
  \qquad
  q_S^\tau(\cdot\mid x,t)
  =\operatorname{softmax}\!\left(z_S(x)_t/\tau\right).
\]
\begin{samepage}
The implemented output-distillation loss is
\begin{equation}
  \mathcal L_{\mathrm{KD}}
  =
  \tau^2
  \frac{
    \sum_{x,t}m_{xt}
    D_{\mathrm{KL}}\!\left(
      p_T^\tau(\cdot\mid x,t)
      \,\|\,
      q_S^\tau(\cdot\mid x,t)
    \right)
  }{
    \sum_{x,t}m_{xt}
  },
  \label{eq:implemented-kd}
\end{equation}
where \(m_{xt}\) is one at non-padding positions and zero otherwise.
Equation~\ref{eq:implemented-kd} averages teacher-to-student token-level \ac{kl} divergence over valid positions.
The factor \(\tau^2\) compensates for the gradient-scale change introduced by temperature softening.
Every executed output-distillation route uses \(\tau=2\), gradient-norm clipping at 1, and AdamW at its library default decay rates and weight decay, with a constant learning rate and no warmup.
Two route types recur below.
An unanchored route trains the student without a teacher and keeps language quality with its own next-token objective, and a teacher-anchored route adds output \ac{kd} from a frozen teacher. \evd{E003}\evd{E005}\evd{E016}
\end{samepage}

The language objective is shared within a paired intervention, but the student initialization and the quality rule that applies differ across families.
Table~\ref{tab:language-training-rules} records those differences.

\begin{table}[tbp]
  \centering
  \footnotesize
  \caption{Language objectives and quality rules by training family.}
  \label{tab:language-training-rules}
  \evd{E003}\evd{E004}\evd{E005}\evd{E008}\evd{E013}\evd{E016}\evd{E017}\evd{E025}\begin{tabularx}{\textwidth}{@{}>{\raggedright\arraybackslash}p{0.18\textwidth}>{\raggedright\arraybackslash}X>{\raggedright\arraybackslash}X>{\raggedright\arraybackslash}X@{}}
    \toprule
    Family & Teacher, student, and initialization & Language objective and corpus & Quality rule \\
    \midrule
    Distillation-headroom analysis & GPT-2 Medium teacher; GPT-2 Small student initialized either from pretrained weights or randomly & Output \ac{kd} on 96,000 WikiText sentences; 2,000 fixed held-out sentences & Held-out perplexity is a convergence measure; the warm and random-start routes differ in initialization and training duration \evd{E003} \\
    \midrule
    Recorded-response intervention (sentence) & Qwen2.5-0.5B student; Qwen2.5-1.5B teacher when the route is teacher-anchored & Masked next-token cross-entropy in the unanchored screen; output \ac{kd} in anchored routes; fixed WikiText quality probe & The recorded-response arm and permuted-response control share the scoring implementation; ordinary \ac{kd} comparisons require perplexity within the tolerance stated in Section~\ref{sec:estimands} \evd{E004}\evd{E005}\evd{E008} \\
    \midrule
    Recorded-response intervention (naturalistic) & Qwen2.5-0.5B student; frozen Qwen2.5-1.5B teacher in anchored routes & Unanchored or output-\ac{kd}-anchored language path, according to the executed grid & The full-parameter route required improved held-out predictivity without a collapse in language quality on the fixed probe; no numeric maximum was fixed in advance \evd{E013}\evd{E017} \\
    \midrule
    Synthetic-response intervention & GPT-2 Medium teacher; pretrained GPT-2 Small student shared within seed & Output \ac{kd} on 95,999 WikiText sentences; 1,999 fixed held-out sentences & Each target arm must satisfy Equation~\ref{eq:quality-eligibility}; saved-student comparisons also apply the within-seed bits-per-byte tolerance of Section~\ref{sec:estimands} \evd{E016}\evd{E025} \\
    \bottomrule
  \end{tabularx}
\end{table}

The headroom family uses seeds 0--2, batch size 16, and maximum length 64.
The pretrained student is trained for one epoch at \(5\times10^{-5}\), whereas the random-start student is trained for two epochs at \(3\times10^{-4}\).
These are separate optimization regimes, not an initialization-only contrast. \evd{E003}
The recorded- and synthetic-response settings are specified with their target paths below.

\subsection{Recorded-response intervention}

For a recorded response \(y_i\), the common sentence-level target path is
\[
  x_i
  \longrightarrow h_{12}(x_i)
  \longrightarrow W h_{12}(x_i)
  \longrightarrow \mathcal L_{\mathrm{target}}.
\]
Here \(h_j\) denotes entry \(j\) in the model's returned hidden-state sequence, whose entry 0 is the embedding output.
The thesis shorthand ``hidden-state index 12'' therefore names the exact attachment returned by the implementation.
The vector \(h_{12}(x_i)\) is mean-pooled over non-padding token positions, and \(W\) is a temporary linear map from the pooled hidden state to the five-\ac{roi} response vector of Appendix~\ref{app:data-targets}, unless the selected loss is readout-free.

\begin{samepage}
The common joint objective is
\begin{equation}
  \mathcal L
  =
  \lambda_{\mathrm{lang}}\mathcal L_{\mathrm{lang}}
  +
  \lambda_{\mathrm{target}}\mathcal L_{\mathrm{target}},
  \qquad
  \lambda_{\mathrm{lang}}=1.
  \label{eq:recorded-objective}
\end{equation}
The first term preserves the declared language objective, while the second exposes the retained student path to the recorded response.
For teacher-anchored routes, \(\mathcal L_{\mathrm{lang}}=\mathcal L_{\mathrm{KD}}\); the unanchored screen instead uses masked next-token cross-entropy.
\end{samepage}

For predictions \(\widehat Y\) and targets \(Y\) with minibatch size \(B\) and target width \(D\), the squared-error target loss is
\[
  \operatorname{MSE}(\widehat Y,Y)
  =
  \frac{1}{BD}
  \sum_{b=1}^{B}
  \sum_{j=1}^{D}
  \left(\widehat Y_{bj}-Y_{bj}\right)^2.
\]
This loss averages equally over minibatch rows and target coordinates.
Consequently, every reported target weight is interpreted under this normalization.

The common sentence-level regime places rank-16 \ac{lora} adapters with scaling 32 and dropout 0.05 in the attention and feed-forward projections of the student family, namely \texttt{c\_attn}, \texttt{c\_proj}, and \texttt{c\_fc} for GPT-2, and \texttt{q\_proj}, \texttt{k\_proj}, \texttt{v\_proj}, \texttt{o\_proj}, \texttt{gate\_proj}, \texttt{up\_proj}, and \texttt{down\_proj} for Qwen2.5 \autocite{hu2022}. \evd{E004}
The pretrained base, input embeddings, and tied output matrix are frozen; the adapters and a co-trained \(W\) are optimizer parameters.
The target gradient reaches \(W\) and the adapters that produce \(h_{12}\), while adapters after the attachment state receive only the language-loss gradient.
After training, the adapters are merged into the student and the temporary target map is discarded.
The merged student is frozen before a fresh evaluation readout is fitted as specified in Appendix~\ref{app:inference}. \evd{E004}\evd{E005}\evd{E008}

Table~\ref{tab:recorded-training-variants} states only the implementation change relative to this common path.

\begingroup
\begin{table}[tbp]
  \footnotesize
  \setlength{\tabcolsep}{3pt}
  \renewcommand{\arraystretch}{1.02}
  \centering
  \caption{Recorded-response intervention variants.}
  \label{tab:recorded-training-variants}
  \begin{tabular}{@{}>{\raggedright\arraybackslash}p{22mm}>{\raggedright\arraybackslash}p{38mm}>{\raggedright\arraybackslash}p{44mm}>{\raggedright\arraybackslash}p{28mm}@{}}
  \toprule
  Variant & Change from the common target path & Executed grid & Retained object \\
  \midrule
  \evd{E004}\evd{E005}\evd{E008}\evd{E011}\evd{E013}\evd{E017}Co-trained \ac{mse} & A new \(W\) is trained jointly with the producing adapters; correctly paired and permuted responses use the same path & Rank 16; seeds 0--2; 3 epochs; \(2\times10^{-4}\); batch 16; target weight 10. Participant-averaged training uses one permutation draw; participant-specific training uses five & Merged student; \(W\) discarded \evd{E004}\evd{E005}\evd{E008} \\
  \midrule
  Frozen readout & \(W_0\) is ridge-fitted on untuned outer-training features with \(\alpha=100\) and remains fixed during intervention training & Same grid and paired-permutation construction as the co-trained \ac{mse} row & Merged student; \(W_0\) discarded \evd{E004} \\
  \midrule
  Contrastive & Symmetric InfoNCE compares normalized \(Wh_{12}\) with the paired five-\ac{roi} response using in-batch negatives \autocite{oord2018} & Rank 16; seeds 0--1; 3 epochs; target weight 1; temperature 0.07; three permutation draws & Merged student; \(W\) discarded \evd{E013} \\
  \midrule
  Increased capacity & The co-trained \ac{mse} path is unchanged; adapter rank increases from 16 to 64 & Seeds 0--1; 6 epochs; target weight 10; five permutation draws; otherwise as in the participant-specific grid & Merged student; \(W\) discarded \evd{E011} \\
  \midrule
  Naturalistic voxelwise & A participant-specific head maps pooled 25-word segments to reliable voxels after a 4-s response delay & UTS01--03; 14 tuning and 6 evaluation stories. The unanchored grid ran all three participants at seeds 0--1, two permutation draws, and target weight 10. The output-\ac{kd}-anchored grid swept target weights 1, 3, and 6 on UTS03 over the 2,801 voxels whose repeat reliability exceeds 0.7 & Merged student; voxel head discarded \evd{E013} \\
  \midrule
  Full parameter & The voxelwise \ac{mse} and output-\ac{kd} anchor update all student parameters except the tied input/output embedding; no \ac{lora} is used & UTS01--03; seeds 0--2; 2 epochs; \(10^{-5}\); batch 16; target weight 10 & Student used for immediate frozen evaluation; voxel head discarded; no reusable student checkpoint retained \evd{E017} \\
  \bottomrule
  \end{tabular}
\end{table}
\endgroup

The unanchored sentence screen also explored alternative target losses, holding the co-trained \ac{mse} grid fixed and changing only \(\mathcal L_{\mathrm{target}}\).
Cosine and squared-Pearson losses use a co-trained target map.
Linear \ac{cka} is readout-free.
It compares the minibatch geometry of \(h_{12}\) and the response matrix directly, so no temporary \(W\) participates in that branch \autocite{kornblith2019}.
These are implementation variants of the recorded-response intervention; their outcomes remain in Section~\ref{sec:results-recorded} and Appendix~\ref{app:secondary}. \evd{E004}

The practical-utility evaluation used this same recorded-response training path, with a Qwen2.5-0.5B student against a Qwen2.5-1.5B teacher at \(h_{12}\), seeds 0--7, target weight 10, 800 training items, three epochs, learning rate \(2\times10^{-4}\), and batch size 16, then 400 items scored from each of the two out-of-domain corpora of Section~\ref{sec:estimands}. \evd{E009}

\subsection{Synthetic-response intervention}

The synthetic-response intervention uses the same joint objective in Equation~\ref{eq:recorded-objective}, but the attachment and trainable path differ.
For GPT-2 input \(x\), the retained student path is
\[
  x
  \longrightarrow e_{\mathrm{tok}}(x)+e_{\mathrm{pos}}(x)
  \longrightarrow h_6(x)
  \longrightarrow h_{7\text{--}12}(x)
  \longrightarrow q_S^\tau(\cdot\mid x),
\]
and the temporary target branch is
\[
  \operatorname{meanpool}\!\left(h_6(x)\right)
  \longrightarrow
  W_{\mathrm{tmp}}\operatorname{meanpool}\!\left(h_6(x)\right)
  \longrightarrow
  \mathcal L_{\mathrm{target}}.
\]
As above, \(h_6\) names hidden-state index 6 in the returned sequence, not an informal count of visible blocks.
The output-\ac{kd} gradient reaches the entire student path.
The target gradient reaches \(W_{\mathrm{tmp}}\), token and positional embeddings, normalization components, and student blocks that produce \(h_6\); blocks after the attachment receive only the output-\ac{kd} gradient.
Because GPT-2 ties its token embedding and output matrix, that retained matrix receives the target gradient through its input role and the output-\ac{kd} gradient through both roles. \evd{E016}

The teacher is frozen, whereas every student parameter and \(W_{\mathrm{tmp}}\) is optimized.
All target arms map the 768-dimensional mean-pooled attachment state to all 20,484 coordinates using the same mean \ac{mse}.
Appendix~\ref{app:data-targets} defines the construction and standardization of the synthetic brain-response target, the text-derived auxiliary target, and the saved block permutation.
The common optimization grid uses 95,999 training sentences, one epoch, learning rate \(5\times10^{-5}\), batch size 4, and maximum length 64. \evd{E016}

\begin{table}[tbp]
  \centering
  \footnotesize
  \caption{Synthetic-response training arms.}
  \label{tab:synthetic-training-arms}
  \evd{E016}\evd{E025}\evd{E026}\begin{tabularx}{\textwidth}{@{}>{\raggedright\arraybackslash}p{0.22\textwidth}>{\raggedright\arraybackslash}X>{\raggedright\arraybackslash}X@{}}
    \toprule
    Arm & Additional supervision & Shared optimization settings \\
    \midrule
    Ordinary \ac{kd} arm & None; \(\lambda_{\mathrm{target}}=0\) and no temporary target map & Pretrained GPT-2 Small initialization, frozen GPT-2 Medium teacher, corpus, output-\ac{kd} loss, and seeds 0--5 \evd{E016}\evd{E025} \\
    \midrule
    Synthetic-response arm & Correctly paired synthetic brain-response target through \(W_{\mathrm{tmp}}\); \(\lambda_{\mathrm{target}}=0.1\) & Same initialization, teacher, corpus, schedule, and seeds 0--5 as the matched ordinary \ac{kd} arm \evd{E016}\evd{E025} \\
    \midrule
    Saved block-permutation sensitivity & Saved block-permuted synthetic target and saved block-permuted text-derived target, each through the same \(W_{\mathrm{tmp}}\) and mean \ac{mse} & Same initialization, teacher, corpus, schedule, target weight, and seeds 0--5 \evd{E016}\evd{E026} \\
    \midrule
    Text-derived auxiliary-control arm & Text-derived auxiliary target through the same target-map shape and mean \ac{mse}; \(\lambda_{\mathrm{target}}=0.1\) & Same pretrained student, teacher, corpus, and schedule; executed for seeds 0--5 \evd{E016}\evd{E026} \\
    \bottomrule
  \end{tabularx}
\end{table}

The text-derived auxiliary-control arm was first run for seeds 0--2 and extended to seeds 3--5 before any participant outcome was scored; all reported analyses use the full six seeds.\evd{E016}

At the final optimizer step, the temporary target map is always discarded.
Whether a run evaluates its student immediately or a later assay reloads a saved checkpoint, every post-training score comes from a readout fitted afresh on the frozen student.
Training can therefore shape the retained student through the target gradient, but the training head itself never computes a target-recovery or biological-transfer score. \evd{E016}\evd{E025}\evd{E030}

The frozen row-twin control is not a training arm.
It is constructed only for the later 50-component linear target-projection assay, and Appendix~\ref{app:data-targets} specifies how it is constructed.
Comparator adequacy and all arm outcomes are reported in Section~\ref{sec:results-synthetic-manipulation} and Appendix~\ref{app:sensitivities}. \evd{E026}\evd{E030}

\subsection{Reproduction scope}

Reproducing a run requires four linked objects.
These are the locked software environment, the executed entry point and arguments, the declared seed and final-state rule, and the retained output.
Table~\ref{tab:training-reproduction} indexes the first three; each experiment's own record states the exact configuration.

\begin{table}[tbp]
  \centering
  \footnotesize
  \caption{Training entry points and retained replay information.}
  \label{tab:training-reproduction}
  \evd{E003}\evd{E004}\evd{E005}\evd{E008}\evd{E011}\evd{E013}\evd{E016}\evd{E017}\evd{E025}\evd{E026}\begin{tabularx}{\textwidth}{@{}>{\raggedright\arraybackslash}p{0.26\textwidth}>{\raggedright\arraybackslash}p{0.30\textwidth}>{\raggedright\arraybackslash}X@{}}
    \toprule
    Family & Executed entry point & Final-state and retained-output rule \\
    \midrule
    Distillation-headroom analysis & \path{scripts/run_kd_alignment.py} & Final warm or random-start student evaluated in memory; the run records its configuration and quality measurements \evd{E003} \\
    \midrule
    Recorded-response intervention (sentence) & \path{scripts/run_brain_lever.py}; target losses in \path{scripts/brain_loss.py} & Final merged student evaluated after each fold--seed cell; temporary target map discarded; no intermediate quality frontier retained \evd{E004}\evd{E005}\evd{E008}\evd{E011} \\
    \midrule
    Recorded-response intervention (naturalistic) & \path{scripts/run_lebel_tune.py} & Final merged or full-parameter student evaluated immediately; temporary voxel head discarded \evd{E013}\evd{E017} \\
    \midrule
    Synthetic-response intervention & \path{scripts/run_tribe_phase3.py} & Final student evaluated in memory; optional checkpoint written only when a save directory is supplied; temporary target map discarded \evd{E016}\evd{E025}\evd{E026} \\
    \bottomrule
  \end{tabularx}
\end{table}

The environment is defined by the Python requirement in \path{pyproject.toml} and the exact package and CUDA build resolved by \path{uv.lock}.
These pin the current environment rather than reconstruct the one every earlier run used.
Each run stores the model identifiers and training configuration it executed.
Many early runs, however, do not record an immutable upstream Hugging Face revision, so exact replay of those runs additionally depends on a retained model cache or a saved student.
This is a real limit on reproduction and is reported rather than filled with an inferred revision.

With the optimization and reproduction scope fixed, Appendix~\ref{app:inference} specifies how the retained students are evaluated and how uncertainty is assigned.
